\section{Background \& Related Work}

\subsection{Evolution through Large Models}
Advancements in language models have enabled new kinds of powerful search algorithms that apply LMs as search operators, e.g.\ to create variation or evaluate solutions. While other search algorithms could also be used, this paper creates a QDAIF algorithm by extending upon Evolution through Large Models (ELM) \citep{lehman2022evolution}, a framework for evolutionary search for code or text that uses LMs to generate intelligent variation (for example through specialized language models trained on code diffs \citep{bradley2023diffmodels}, or through simple few-shot prompting \citep{meyerson2023language,chen2023evoprompting}). Most QDAIF results in this paper generate new search candidates through Language Model Crossover (LMX) \citep{meyerson2023language}, a recent and general few-shot prompting approach that can evolve e.g. mathematical expressions, sentences, Python programs, and prompts for text-to-image models, by leveraging in-context learning capabilities of LMs \citep{brown2020language}. The approach is simple: A few existing search candidates are concatenated into a prompt, predisposing the LM to generate new, similar candidates. In this way, LMX enables creating intelligent variation without requiring any specially-trained models. Our experimental implementation builds on OpenELM \citep{Bradley_OpenELM_2023}, a versatile open-source Python library designed for research into LM-based evolutionary algorithms.

\subsection{Quality Diversity Algorithms}\label{main:qd_background}
Traditional optimization algorithms aim to discover a single high-quality solution, which while appropriate for many situations, can fail to \emph{illuminate} the full range of possible high-quality solutions. For creative and design problems in particular, a user may want to choose what they think is most appropriate from a diversity of such candidates. In contrast, Quality Diversity (QD) algorithms aim to optimize not just for a single optimal solution, but for a diverse set of high-quality solutions \citep{lehman2011evolving, mouret2015illuminating, pugh2016quality, fontaine2021differentiable}. QD algorithms can thus provide a richer landscape of solutions, enabling adaptability and flexibility in addressing multifaceted challenges \citep{cully2015robots}.  In addition to a quality measure (i.e. an objective function), QD requires a metric such that it can encourage desired axes of diversity. For instance, \citet{lehman2022evolution} evolved Python programs to design varied locomoting robots, where the diversity dimensions are the robot's height, width, and mass.

A significant limitation of existing QD algorithms lies in their reliance on low-level quality and diversity measures \citep{mouret2015illuminating}. This requirement confounds applying QD algorithms to complex and creative domains, such as the creative writing ones explored in this paper. Intuitively, such measures (e.g. sensor readings \citep{cully2015robots}, feature engineering \citep{manning2009introduction}) lack the subtlety and depth needed to capture the complexities of human creativity and intuition, e.g.\ nuances, moods, or cultural references that resonate in human experience. Interestingly, from having trained on vast amounts of human-generated data, LMs can begin to emulate such human-nuanced judgments (cf. \cref{main:ai_feedback}). Thus, by employing an LM to evaluate both quality and diversity, QDAIF significantly simplifies and enlarges the range of domains QD can be applied to. 

Feedback from learned ML models has been used in prior work to reduce the need for hand-crafted heuristics or expensive ground-truth evaluations. In model-based QD, learned feedback is supplied by surrogate models. \citet{gaier2017data} introduced the use of surrogates (via a Gaussian process) to predict fitness (quality). Subsequently, \citet{keller2020model} introduced a learned model to predict both fitness and behavior characteristics (diversity), becoming a standard approach \citep{lim2021dynamics,lim2022learning,zhang2022deep,bhatt2022deep}. Surrogate models require domain-specific training data to update their predictions on a limited domain, whereas AI feedback leverages off-the-shelf instruction-tuned LMs \citep{chung2022scaling,ouyang2022training} to automate expensive human feedback for a variety of evaluation tasks. More recently, \citet{fontaine2021differentiable} utilized CLIP embeddings \citep{radford2021learning} as both quality and diversity measures to navigate the search space of StyleGAN \citep{karras2019style}, producing a range of faces with the desired characteristic (e.g.~"A person with red hair"). We show that using pre-trained surrogate models is more prone to reward hacking in the natural language case \citep{skalse2022defining,lehman2019surprising} (cf.~\cref{main:alt_feedback}). Hence, QDAIF capitalizes on the strengths of general-purpose LMs for evaluating generated solutions.

\subsection{AI Feedback}\label{main:ai_feedback}
Recent months have seen a surge in research that leverages LMs to provide feedback on the training, evaluation, or problem-solving capabilities of other LMs \citep{bai2022constitutional, perez2022discovering, shinn2023reflexion, wang2023voyager, colas2023augmenting, zhang2023omni, lee2023rlaif}. \citet{bai2022constitutional} show that using LM-generated critiques and refinements has been instrumental in enhancing performance on metrics like helpfulness and harmlessness. One particularly promising direction for AI feedback is self-refinement, where LMs evaluate and score their own generations, and then iteratively improve their output \citep{bai2022constitutional,madaan2023self}. Self-refinement has demonstrated significant improvement in output quality as gauged by human evaluators \citep{madaan2023self}, underscoring generation-discrimination discrepancy \citep[p.12]{saunders2022self}, meaning that it is often easier for a model to evaluate the quality of a generation than to generate the same high-quality text. Complementary to single-objective optimization with self-refine, QDAIF utilizes AI feedback to assess diversity in addition to quality, facilitating more varied and improved text generation over multiple iterations of refinement through evolution.





